\documentclass[11pt]{article}

\usepackage[T1]{fontenc}
\usepackage[utf8]{inputenc}
\usepackage{lmodern}
\usepackage{amsmath,amssymb,amsthm,mathtools}
\usepackage{geometry}
\usepackage{booktabs}
\usepackage{array}
\usepackage{longtable}
\usepackage{hyperref}
\usepackage{enumitem}

\geometry{margin=1in}
\setlist[itemize]{leftmargin=1.5em}
\setlist[enumerate]{leftmargin=1.75em}
\emergencystretch=3em

\newcommand{\code}[1]{\texttt{\detokenize{#1}}}
\newcommand{\Fq}{\mathbf{F}_q}
\newcommand{\Z}{\mathbb{Z}}
\newcommand{\Rq}{\mathbb{Z}/q\mathbb{Z}}
\newcommand{\Range}[2]{\{#1,\ldots,#2\}}
\newcommand{\B}{\mathsf{B}}
\newcommand{\W}{\mathsf{W32}}
\newcommand{\bw}{\mathsf{bw32}}
\newcommand{\repr}{\mathsf{repr}}
\newcommand{\poly}{\mathsf{poly}}
\newcommand{\mont}{\mathsf{mont}}
\newcommand{\Proof}{\par\noindent\emph{Proof.} }
\newcommand{\Lemma}[1]{\paragraph{\code{#1}.}}

\title{Mathematical Notes for the Lemmas in \texorpdfstring{\code{RefJasminNTT.ec}}{RefJasminNTT.ec}}
\author{Codex Proof Notes}
\date{28 April 2026}

\begin{document}
\maketitle

\begin{abstract}
These notes give a mathematician-style account of every lemma and equivalence
proved in \code{RefJasminNTT.ec}.  The EasyCrypt file is a refinement layer:
it proves that the extracted Jasmin forward and inverse NTT cores refine the
imperative reference procedures in \code{NTT_Fq.NTT}.  The central mathematical
themes are representation of signed machine words as field elements, Montgomery
table cancellation, pointwise word-size bounds, and loop-level pRHL invariants.
The presentation below follows the order of the EasyCrypt file and explains
what each lemma proves and why the proof is valid.
\end{abstract}

\tableofcontents

\section{Notation and Proof Obligations}

All coefficients live in the finite field \(\Fq=\Rq\), where the HAETAE
modulus is \(q=64513\).  A 32-bit machine word \(w\) is interpreted as the
field element
\[
  \operatorname{word}(w)=\operatorname{incoeff}(\W.\operatorname{to\_sint}(w)).
\]
For a byte array \(rp:\code{BArray1024.t}\), the corresponding coefficient
array is obtained pointwise by reading the \(i\)-th 32-bit lane:
\[
  \operatorname{poly}(rp)[i]
  =
  \operatorname{word}(\code{BArray1024.get32}\;rp\;i).
\]
The relation \(\code{NTT_Fq.poly_repr}(rp,p)\) states
\(\operatorname{poly}(rp)=p\).  The bounded relation
\(\code{NTT_Fq.poly_repr_bound}(rp,p,s)\) additionally states that every lane
of \(rp\) satisfies the signed \(s\)-bit word bound \(\bw(-,s)\).

The proof file also uses a pointwise bound function \(bd:\Z\to\Z\):
\[
  \code{barray256_bound_by}(rp,bd)
  \quad\Longleftrightarrow\quad
  \forall i\in\Range{0}{255}.\,
    \bw(\code{BArray1024.get32}\;rp\;i,bd(i)).
\]
This permits a loop invariant to say that already processed coefficients have
one more bit of headroom than untouched coefficients.

\section{Integer Shift Lemmas}

\Lemma{int_shr1_div2}
For every non-negative integer \(x\), the EasyCrypt integer operation
\(x\mathbin{\code{`|>>}}1\) equals Euclidean division \(x\%/2\).
\Proof
The operator \(\code{`|>>}\) is the integer right-shift notation expanded by
EasyCrypt into the corresponding arithmetic division by a power of two.  With
shift amount \(1\), that power is \(2\).  Since \(x\ge 0\), Euclidean division
and the intended non-negative shift semantics coincide.

\Lemma{int_shl1_mul2}
For every integer \(x\), the EasyCrypt integer operation
\(x\mathbin{\code{`<<}}1\) equals \(2x\).
\Proof
Expanding the definition of \(\code{`<<}\), left shift by one is multiplication
by \(2^1\).  Hence the result is \(x\cdot 2\).

\section{Machine-Word Addition and Subtraction}

\Lemma{word_to_coeff_add}
Assume two words \(a,b\) satisfy bounds \(\bw(a,a_s)\) and \(\bw(b,b_s)\), with
\(0\le a_s,b_s<31\).  Then
\[
  \operatorname{word}(a+b)
  =
  \operatorname{word}(a)+\operatorname{word}(b).
\]
\Proof
The hypotheses imply that the signed representatives of \(a\) and \(b\) lie in
intervals contained in \((-2^{30},2^{30})\).  Therefore their sum lies inside
the range where \(\W.\operatorname{to\_sint}\) of the word addition is the
integer sum of signed representatives.  EasyCrypt's lemma
\code{W32.to_sintD_small} gives this equality at the integer level.  Applying
\(\operatorname{incoeff}\) and using additivity of \(\operatorname{incoeff}\)
gives the field equality.

\Lemma{word_to_coeff_sub}
Under the same hypotheses as the preceding lemma,
\[
  \operatorname{word}(a-b)
  =
  \operatorname{word}(a)-\operatorname{word}(b).
\]
\Proof
The signed bounds again put both operands in a range where word subtraction
does not wrap with respect to signed interpretation.  The lemma
\code{W32.to_sintB_small} identifies the signed representative of \(a-b\) with
the integer difference.  The conclusion follows from the compatibility of
\(\operatorname{incoeff}\) with subtraction.

\Lemma{forward_butterfly_bounds}
If \(s\) and \(t\) satisfy bounds \(a_s\) and \(t_s\), with both exponents in
\([0,30]\), then both butterfly outputs \(s+t\) and \(s-t\) satisfy bound
\(\max(a_s,t_s)+1\).
\Proof
The signed magnitude of each operand is bounded by the appropriate power of
two.  The sum and difference are therefore bounded by twice the larger bound,
which is exactly one additional bit.  The EasyCrypt proof invokes the word
arithmetic lemmas \code{Fq.add_corr} and \code{Fq.sub_corr}.

\Lemma{inverse_butterfly_bounds}
If \(t_0\) and \(coeff\) satisfy bounds \(a_s\) and \(b_s\), then both
\(t_0+coeff\) and \(t_0-coeff\) satisfy bound \(\max(a_s,b_s)+1\).
\Proof
This is the same signed-interval argument as for
\code{forward_butterfly_bounds}, specialized to the inverse butterfly
temporaries.  Addition and subtraction each cost at most one bit.

\section{Montgomery Constant and Table Cancellation}

\Lemma{unit_R_ref}
The Montgomery constant \(R=\code{NTT_Fq.R}\) is a unit of \(\Fq\).
\Proof
The constant \(R\) is the residue of the Montgomery radix modulo \(q\).  The
proof expands \(R\), reduces the claim to a non-zero residue modulo \(q\), and
uses the unit criterion for the field \(\Fq\).

\Lemma{inv_R_ref}
The inverse of \(R\) in \(\Fq\) is \(\operatorname{incoeff}(50386)\).
\Proof
Since \(R\) is a unit by \code{unit_R_ref}, it suffices to multiply both sides
by \(R\).  Expanding \(R\) and reducing modulo \(q\), EasyCrypt checks
\[
  R\cdot 50386 \equiv 1 \pmod q.
\]
The field equality follows from the injective residue equality used by
\(\operatorname{incoeff}\).

\Lemma{array256_mont_get}
For \(i\in\Range{0}{255}\),
\[
  \code{NTT_Fq.array256_mont}(p)[i]=p[i]\cdot R.
\]
\Proof
The operator \code{array256_mont} is an \code{Array256.map}/\code{init}
definition that multiplies every coefficient by \(R\).  Expanding the
definition and evaluating the initializer at an in-range index proves the
claim.

\Lemma{mont_cancel}
For \(i\in\Range{0}{255}\),
\[
  \code{NTT_Fq.array256_mont}(p)[i]\cdot R^{-1}=p[i].
\]
\Proof
By \code{array256_mont_get}, the left-hand side is
\(p[i]\cdot R\cdot R^{-1}\).  The lemma \code{inv_R_ref} identifies \(R^{-1}\),
and the unit law \(R R^{-1}=1\) reduces the expression to \(p[i]\).

\Lemma{jzetas_get_cancel}
For \(1\le i<256\),
\[
  \operatorname{word}(\code{jzetas}[i])\cdot R^{-1}
  =
  \code{NTT_Fq.zetas}[i].
\]
\Proof
The table lemma \code{NTT_Fq.jzetas_get} states that the extracted table entry
at index \(i\) is the Montgomery encoding of the abstract forward twiddle:
\(\operatorname{word}(\code{jzetas}[i])=(\code{array256_mont zetas})[i]\).
Applying \code{mont_cancel} cancels the factor \(R\).

\Lemma{array256_mont_inv_get}
For \(i\in\Range{0}{255}\),
\[
  \code{NTT_Fq.array256_mont_inv}(p)[i]=p[i]\cdot R.
\]
\Proof
This is the inverse-table analogue of \code{array256_mont_get}.  It follows by
expanding \code{array256_mont_inv} as the pointwise multiplication map.

\Lemma{mont_inv_cancel}
For \(i\in\Range{0}{255}\),
\[
  \code{NTT_Fq.array256_mont_inv}(p)[i]\cdot R^{-1}=p[i].
\]
\Proof
Combine \code{array256_mont_inv_get} with the identity \(R R^{-1}=1\), using
\code{inv_R_ref}.

\Lemma{jzetas_inv_get_cancel}
For \(i\in\Range{0}{254}\),
\[
  \operatorname{word}(\code{jzetas_inv}[i])\cdot R^{-1}
  =
  \code{NTT_Fq.zetas_inv}[i].
\]
\Proof
The imported lemma \code{NTT_Fq.jzetas_inv_get} identifies the concrete
inverse table entry with the Montgomery encoding of \(\code{zetas_inv}[i]\).
The cancellation lemma \code{mont_inv_cancel} removes \(R\).  The range stops
at \(254\) because index \(255\) is the special final scaling constant.

\Lemma{jzetas_inv_255_scale_cancel}
At the final inverse-table index,
\[
  \operatorname{word}(\code{jzetas_inv}[255])\cdot R^{-1}
  =
  \code{NTT_Fq.scale255}\cdot R.
\]
\Proof
This is exactly the table fact \code{NTT_Fq.jzetas_inv_255_scaleE}.  It states
that the extracted word at index \(255\) is chosen so that, after one
Montgomery multiplication, the effect is multiplication by
\(\code{scale255}\) and retention of one Montgomery factor.

\section{Montgomery Multiplication Lemmas}

\Lemma{fqmul_bw32_16}
If the integer product \(aa\cdot bb\) lies in the signed Montgomery reducer
range
\[
 -\frac{R}{2}q\le aa\,bb<\frac{R}{2}q,
\]
then a call to \code{Hpoly_extract.M.__fqmul} with operands whose signed values
are \(aa\) and \(bb\) returns a 16-bit bounded word.
\Proof
The extracted multiplication routine is related to the signed Montgomery
reducer by \code{Fq.fqmul_corr_h}.  The reducer correctness lemma
\code{Fq.SignedReductions.SREDCp_corr} says the result lies in
\((-q,q)\).  Since \(q<2^{16}\), the returned signed word satisfies
\(\bw(-,16)\).

\Lemma{fqmul_word_to_coeff_mul}
Under the same product-range hypothesis, \code{__fqmul} returns a word whose
field interpretation is
\[
  \operatorname{incoeff}(aa)\operatorname{incoeff}(bb)R^{-1}.
\]
\Proof
Again use \code{Fq.fqmul_corr_h} to replace the procedure result by the signed
Montgomery reduction of \(aa\,bb\).  The congruence part of
\code{SREDCp_corr} states that this reduction is congruent to
\(aa\,bb\,R^{-1}\) modulo \(q\).  Translating through
\(\operatorname{incoeff}\) and \code{inv_R_ref} gives the field equality.

\Lemma{fqmul_word_to_coeff_mul_h}
This is the Hoare-form version of \code{fqmul_word_to_coeff_mul}, with the
product-range hypothesis placed in the procedure precondition.
\Proof
The proof is the same congruence argument as above.  The only difference is
logical: the product bound is destructed from the Hoare precondition rather
than supplied as a lemma parameter.

\Lemma{fqmul_word_to_coeff_mul_bound_h}
This combines the semantic conclusion of \code{fqmul_word_to_coeff_mul_h} with
the 16-bit output bound of \code{fqmul_bw32_16}.
\Proof
Use Hoare consequence twice.  The first consequence supplies the field-value
postcondition from \code{fqmul_word_to_coeff_mul_h}; the second invokes
\code{Fq.fqmul_corr_h} and \code{SREDCp_corr} to obtain the 16-bit bound.

\Lemma{fqmul_word_to_coeff_mul_bound_ph}
The probabilistic Hoare version of the previous lemma holds with probability
one.
\Proof
The routine \code{__fqmul} is lossless by \code{Fq.fqmul_ll}.  Combining
losslessness with the ordinary Hoare triple
\code{fqmul_word_to_coeff_mul_bound_h} yields the probability-one
\code{phoare} statement.

\section{Pointwise Bound Infrastructure}

The file defines
\[
\code{barray256_bound_by}(rp,bd)
\]
and
\[
\code{poly_repr_bound_by}(rp,p,bd).
\]
The first is a pointwise signed-word bound for all 256 lanes; the second adds
the field representation relation.  The following lemmas are the elementary
logic of these definitions.

\Lemma{barray256_bound_by_get}
If \(\code{barray256_bound_by}(rp,bd)\) holds and \(i\in\Range{0}{255}\), then
the \(i\)-th lane satisfies \(\bw(\code{get32}\;rp\;i,bd(i))\).
\Proof
This is immediate from universal elimination in the definition of
\code{barray256_bound_by}.

\Lemma{barray256_bound_by_weaken}
If \(bd_1(i)\le bd_2(i)\) for every in-range \(i\), then a buffer bounded by
\(bd_1\) is also bounded by \(bd_2\).
\Proof
At each index, use \code{NTT_Fq.bw32_weaken}, which says that a signed word
bounded by a smaller exponent is bounded by a larger exponent.

\Lemma{barray256_bound_by_set32}
If \(rp\) is bounded by \(bd\), \(i\) is in range, and the new word \(w\)
satisfies \(\bw(w,bd(i))\), then \(\code{set32}\;rp\;i\;w\) is still bounded by
\(bd\).
\Proof
For an arbitrary lane \(j\), use the byte-array get-after-set lemma:
the read is \(w\) when \(j=i\), and is the old lane otherwise.  The case
\(j=i\) uses the hypothesis on \(w\); the other case uses the old bound.

\Lemma{barray256_bound_by_set32_change}
Suppose \(rp\) is bounded by \(bd\).  If a store at \(i\) writes a word bounded
by \(bd'(i)\), and every untouched lane has \(bd(k)\le bd'(k)\), then the
updated buffer is bounded by \(bd'\).
\Proof
The same get-after-set case split is used.  At the updated index, the new word
has the desired bound by hypothesis.  At all other indices, the old bound is
weakened from \(bd(k)\) to \(bd'(k)\).

\Lemma{poly_repr_bound_by_repr}
The relation \(\code{poly_repr_bound_by}(rp,p,bd)\) implies
\(\code{NTT_Fq.poly_repr}(rp,p)\).
\Proof
This is the first conjunct of the definition.

\Lemma{poly_repr_bound_by_bound}
The relation \(\code{poly_repr_bound_by}(rp,p,bd)\) implies
\(\code{barray256_bound_by}(rp,bd)\).
\Proof
This is the second conjunct of the definition.

\Lemma{fqmul_product_bound_16_24}
If \(a\) is 16-bit bounded and \(b\) is 24-bit bounded, then
\[
 -\frac{R}{2}q
 \le
 \W.\operatorname{to\_sint}(a)\W.\operatorname{to\_sint}(b)
 <
 \frac{R}{2}q.
\]
\Proof
The signed bounds give
\(|\W.\operatorname{to\_sint}(a)|<2^{16}\) and
\(|\W.\operatorname{to\_sint}(b)|<2^{24}\).  With the concrete HAETAE values
of \(R\) and \(q\), the product interval is strictly contained in the
Montgomery reducer input interval.  EasyCrypt discharges the remaining
integer arithmetic by normalization and SMT.

\Lemma{jzetas_bound16}
Every forward twiddle table entry \(\code{jzetas}[i]\), for \(1\le i<256\), is
16-bit bounded.
\Proof
Expand the concrete list \(\code{Hpoly_extract.jzetas}\) and use
\code{BArray1024.get32_of_list32}.  Each of the 255 in-range constants is
checked by reducing the signed word conversion and verifying its magnitude is
below \(2^{16}\).

\Lemma{jzetas_inv_bound16}
Every inverse twiddle table entry \(\code{jzetas_inv}[i]\), for
\(0\le i<256\), is 16-bit bounded.
\Proof
The proof is the same finite table enumeration as \code{jzetas_bound16}, now
over all 256 entries of the inverse table, including the final scaling entry.

\Lemma{barray256_bound_by_const}
A uniform bound \(\code{NTT_Fq.barray256_bound}(rp,s)\) implies the pointwise
bound by the constant function \(i\mapsto s\).
\Proof
Both definitions quantify over the same 256 indices.  Expanding
\code{const_bound} makes the conclusions identical.

\Lemma{barray256_bound_by_const_bound}
Conversely, a pointwise bound by the constant function \(i\mapsto s\) implies
\(\code{NTT_Fq.barray256_bound}(rp,s)\).
\Proof
Expand \code{const_bound}; the pointwise and uniform definitions coincide.

\section{Forward Stage Arithmetic}

The forward stage definitions encode the descending layer sequence:
\[
128,64,32,16,8,4,2,1,0.
\]
The associated bound table is
\[
\begin{array}{c|ccccccccc}
\ell&128&64&32&16&8&4&2&1&0\\
\midrule
\code{fwd_stage_bound}(\ell)&16&17&18&19&20&21&22&23&24
\end{array}
\]
and the twiddle base table is
\[
\begin{array}{c|ccccccccc}
\ell&128&64&32&16&8&4&2&1&0\\
\midrule
\code{fwd_zbase}(\ell)&0&1&3&7&15&31&63&127&255.
\end{array}
\]

\Lemma{fwd_stage_bound_range}
For every legal forward length, \(\code{fwd_stage_bound}(\ell)\) lies between
0 and 24.
\Proof
The predicate \code{fwd_len_ok} is a finite disjunction over the nine legal
values.  Checking the definition of \code{fwd_stage_bound} on each case gives
one of \(16,\ldots,24\).

\Lemma{fwd_stage_bound_pos}
For every legal forward length, \(\code{fwd_stage_bound}(\ell)\ge 0\).
\Proof
This is the lower half of \code{fwd_stage_bound_range}.

\Lemma{fwd_stage_bound_lt31}
For every legal forward length, \(\code{fwd_stage_bound}(\ell)<31\).
\Proof
The upper half of \code{fwd_stage_bound_range} gives at most \(24\), hence
strictly less than \(31\).

\Lemma{fwd_stage_bound_ge16}
For every legal forward length, \(\code{fwd_stage_bound}(\ell)\ge 16\).
\Proof
Case analysis on the finite definition gives values \(16,\ldots,24\).

\Lemma{fwd_stage_next}
If \(\ell\) is a positive legal forward length, then
\[
  \code{fwd_stage_bound}(\ell/2)
  =
  \code{fwd_stage_bound}(\ell)+1.
\]
\Proof
The positive legal lengths are \(128,\ldots,1\).  Halving moves one step to
the right in the bound table, increasing the bound by one.

\Lemma{fwd_len_next_ok}
If \(\ell\) is a positive legal forward length, then \(\ell/2\) is again a
legal forward length.
\Proof
The allowed positive lengths are powers of two from \(128\) to \(1\).  Dividing
by two gives the next entry in the list, with \(1/2=0\) in integer division,
which is included as the loop-exit length.

\Lemma{fwd_zbase_next}
For a positive legal length,
\[
  \code{fwd_zbase}(\ell/2)
  =
  \code{fwd_zbase}(\ell)+256/(2\ell).
\]
\Proof
Each forward layer consumes \(256/(2\ell)\) twiddle entries.  The explicit
\code{fwd_zbase} table records the cumulative number of entries consumed
before the layer.  Case analysis over \(\ell=128,\ldots,1\) verifies the
identity.

\Lemma{fwd_zbase_range}
For every legal forward length, \(\code{fwd_zbase}(\ell)\) lies between 0 and
255.
\Proof
Immediate from the explicit table \(0,1,3,7,15,31,63,127,255\).

\Lemma{fwd_read_range}
Assume a positive legal length, a block start \(0\le start<256\), and
\[
start=2\ell\,(zc-\code{fwd_zbase}(\ell)).
\]
Then the forward table read index \(zc+1\) satisfies \(1\le zc+1<256\).
\Proof
The equation expresses \(zc-\code{fwd_zbase}(\ell)\) as the number of completed
blocks in the layer.  Since \(start\) is a valid block start strictly below
256, this number lies in the range of blocks of the current layer.  Combining
this with the explicit \code{fwd_zbase} table gives \(zc+1\in\Range{1}{255}\).

\Lemma{fwd_block_end_bound}
Under the hypotheses of \code{fwd_read_range},
\[
  start+2\ell\le 256.
\]
\Proof
The same block-count equation says that \(start\) is the beginning of a full
block of size \(2\ell\) in a 256-entry array.  Legal layer sizes divide 256,
so the end of the block is still within the array.

\Lemma{fwd_stage_zbase_exit}
If a positive forward layer exits the middle loop, so that
\[
  256=2\ell\,(zc-\code{fwd_zbase}(\ell)),
\]
then \(zc=\code{fwd_zbase}(\ell/2)\).
\Proof
At layer exit, the number of consumed blocks is \(256/(2\ell)\).  Adding this
quantity to the old zeta base gives the next zeta base by
\code{fwd_zbase_next}.

\Lemma{fwd_middle_to_inner}
The middle-loop bound at block start equals the inner-loop bound with
\(j=start\):
\[
\code{fwd_middle_bound}(s,start,i)
=
\code{fwd_inner_bound}(s,\ell,start,start,i).
\]
\Proof
At the beginning of a block, no coefficient of that block has been processed
yet.  Expanding both definitions shows that the only entries assigned the
larger bound are exactly those with \(0\le i<start\).

\Lemma{fwd_inner_current}
For \(j\) inside the current block, the current butterfly inputs \(j\) and
\(j+\ell\) have bound \(s\), and after advancing to \(j+1\) both positions
have bound \(s+1\).
\Proof
Unfold \code{fwd_inner_bound}.  Before processing the pair, neither \(j\) nor
\(j+\ell\) lies in the completed subranges.  After incrementing \(j\), both
positions lie in completed subranges.

\Lemma{fwd_inner_weaken_unchanged}
For any index \(k\) distinct from \(j\) and \(j+\ell\), the inner bound at
\(j\) is no larger than the inner bound at \(j+1\).
\Proof
Advancing the inner loop only adds the two current positions to the processed
sets.  All other positions either keep the same status or move from smaller to
larger bound.  The result follows by unfolding the range predicates.

\Lemma{fwd_inner_exit_to_middle}
At \(j=start+\ell\), the inner-loop bound equals the middle-loop bound for the
next block start \(start+2\ell\).
\Proof
At inner-loop exit, both halves of the block have been processed.  The union of
the old completed prefix, the first half of the block, and the second half of
the block is precisely the prefix \(0\le i<start+2\ell\).

\Lemma{fwd_middle_exit_const}
If \(start\ge 256\), then every in-range index has
\[
\code{fwd_middle_bound}(s,start,i)=s+1.
\]
\Proof
Every \(i\in\Range{0}{255}\) satisfies \(i<start\), so the first branch of
\code{fwd_middle_bound} applies.

\Lemma{fwd_middle_start_const}
At \(start=0\),
\[
\code{fwd_middle_bound}(s,0,i)=s.
\]
\Proof
No index satisfies \(0\le i<0\), so the default branch of the definition
applies.

\section{Forward Butterfly and Twiddle Lemmas}

\Lemma{forward_twiddle_product}
Given a representation relation \(rp\sim p\), a valid forward twiddle index
\(zc\), and operands \(ai=\code{jzetas}[zc]\) and
\(bi=rp[j+\ell]\), a call to \code{__fqmul} returns a word representing
\[
  \code{zetas}[zc]\cdot p[j+\ell].
\]
\Proof
The representation relation gives
\[
p[j+\ell]=\operatorname{word}(bi).
\]
The multiplication correctness lemma gives
\[
\operatorname{word}(res)
=\operatorname{word}(ai)\operatorname{word}(bi)R^{-1}.
\]
By \code{jzetas_get_cancel},
\(\operatorname{word}(ai)R^{-1}=\code{zetas}[zc]\).  Associativity and
commutativity in \(\Fq\) yield the result.

\Lemma{forward_twiddle_product_bound_h}
This strengthens \code{forward_twiddle_product} with the postcondition
\(\bw(res,16)\).
\Proof
The semantic part is \code{forward_twiddle_product}.  The bound part is
\code{fqmul_bw32_16}, using the same product-range hypothesis.

\Lemma{forward_twiddle_product_bound_ph}
This is the probability-one \code{phoare} form of
\code{forward_twiddle_product_bound_h}.
\Proof
Combine losslessness of \code{__fqmul} with the ordinary Hoare triple.

\Lemma{forward_twiddle_product_bound_call_h}
This is a call-friendly Hoare triple for \code{__fqmul}: all equalities tying
local variables to the table entry and coefficient word are included in the
precondition, and the postcondition gives both the abstract twiddle product and
the 16-bit result bound.
\Proof
Apply \code{fqmul_word_to_coeff_mul_bound_h}.  Then rewrite the local-variable
equalities, use \code{poly_repr_get} for \(p[j+\ell]\), and use
\code{jzetas_get_cancel} to remove the Montgomery factor.

\Lemma{forward_twiddle_product_bound_call_ph}
The probability-one version of \code{forward_twiddle_product_bound_call_h}.
\Proof
Use \code{Fq.fqmul_ll} and the Hoare triple
\code{forward_twiddle_product_bound_call_h}.

\Lemma{forward_twiddle_value}
If a word \(t\) is known to satisfy
\[
\operatorname{word}(t)=\operatorname{word}(ai)\operatorname{word}(bi)R^{-1},
\]
with \(ai=\code{jzetas}[zc]\) and \(bi=rp[j+\ell]\), then
\[
\operatorname{word}(t)=\code{zetas}[zc]p[j+\ell].
\]
\Proof
This is the pure post-call version of \code{forward_twiddle_product}.  The
representation relation identifies \(bi\) with \(p[j+\ell]\), and
\code{jzetas_get_cancel} identifies the table entry after multiplying by
\(R^{-1}\).

\Lemma{forward_butterfly_store}
If \(rp\sim p\), \(s=rp[j]\), and the word operations \(s-t\) and \(s+t\) are
within their signed ranges, then the two stores
\[
rp[j+\ell]\leftarrow s-t,\qquad rp[j]\leftarrow s+t
\]
represent the polynomial update
\[
p[j+\ell]\leftarrow p[j]-\operatorname{word}(t),\qquad
p[j]\leftarrow p[j]+\operatorname{word}(t).
\]
\Proof
The representation relation gives \(p[j]=\operatorname{word}(s)\).  The lemmas
\code{word_to_coeff_sub} and \code{word_to_coeff_add} convert the stored words
to the corresponding field subtraction and addition.  Two applications of
\code{NTT_Fq.poly_repr_set32} complete the representation proof.

\Lemma{forward_butterfly_store_bound}
Under a uniform \(s\)-bit input bound and a 16-bit bound on \(t\), the two
stores above produce a representation bounded by \(\max(s,16)+1\).
\Proof
The representation part is \code{forward_butterfly_store}.  The bound part
uses \code{forward_butterfly_bounds} for the two updated words, weakens all
untouched lanes to the larger bound, and applies the uniform
\code{barray256_bound_set32} lemmas.

\Lemma{forward_butterfly_store_bound_by}
This is the pointwise-bound version of \code{forward_butterfly_store_bound}.
It assumes that the current pair \(j,j+\ell\) has bound \(a_s\), and that the
target bound \(bd'\) raises exactly those two positions to \(a_s+1\).
\Proof
Use \code{forward_butterfly_store} for the representation.  Use
\code{forward_butterfly_bounds} for the two newly written words.  The first
store is handled by \code{barray256_bound_by_set32_change}; the second by
\code{barray256_bound_by_set32}.  The hypothesis on untouched lanes supplies
all weakenings.

\Lemma{forward_spec_updateE}
The two syntactic forms of the forward specification update are extensionally
equal:
\[
(p[j+\ell\leftarrow p[j]-t])[j\leftarrow(p[j+\ell\leftarrow p[j]-t])[j]+t]
=
(p[j+\ell\leftarrow p[j]-t])[j\leftarrow p[j]+t].
\]
\Proof
Since \(\ell>0\), the indices \(j\) and \(j+\ell\) are distinct.  Extensional
array equality reduces the claim to an arbitrary index \(i\); the get-after-set
rules split on whether \(i\) is \(j\), \(j+\ell\), or neither.

\Lemma{forward_butterfly_step_bound_by}
One full forward butterfly step preserves representation and updates the
pointwise inner-loop bound from
\(\code{fwd_inner_bound}(s,\ell,start,j)\) to
\(\code{fwd_inner_bound}(s,\ell,start,j+1)\).
\Proof
\code{fwd_inner_current} shows that the two current inputs have bound \(s\).
The 16-bit multiplication result is the \(t\) input to the store.  Apply
\code{forward_butterfly_store_bound_by}; the untouched-lane weakening is
exactly \code{fwd_inner_weaken_unchanged}.  Finally
\code{forward_spec_updateE} aligns the array update with the right-hand
imperative specification.

\Lemma{forward_twiddle_product_bound_from_inner}
Inside the forward inner-loop invariant, the product
\[
\code{jzetas}[zc]\cdot rp[j+\ell]
\]
meets the signed Montgomery reducer input range.
\Proof
\code{jzetas_bound16} gives the 16-bit twiddle bound.  By
\code{fwd_inner_current}, the current second butterfly input has bound
\code{fwd_stage_bound} at length \(\ell\); this bound is at most 24 by
\code{fwd_stage_bound_range}.  Weakening gives a 24-bit coefficient bound.
Now apply \code{fqmul_product_bound_16_24}.

\Lemma{forward_butterfly_repr}
If \(t\) already represents the abstract product
\(\code{zetas}[zc]p[j+\ell]\), then the two Jasmin stores represent the forward
specification butterfly update.
\Proof
Apply \code{forward_butterfly_store} and rewrite
\(\operatorname{word}(t)\) to the abstract product.

\Lemma{forward_butterfly_step}
If \(t\) satisfies the Montgomery product equation using the concrete forward
twiddle and input word, then the two stores represent the abstract forward
butterfly update.
\Proof
First use \code{forward_twiddle_value} to convert the Montgomery product
equation into the abstract product equation.  Then apply
\code{forward_butterfly_repr}.

\section{Inverse Stage Arithmetic}

The inverse stage definitions encode the ascending layer sequence:
\[
1,2,4,8,16,32,64,128,256.
\]
The predicate \code{inv_len_ok} is the finite disjunction recognizing exactly
these legal loop lengths.
The bound table is again \(16,\ldots,24\), while the inverse zeta base table is
\[
\begin{array}{c|ccccccccc}
\ell&1&2&4&8&16&32&64&128&256\\
\midrule
\code{inv_zbase}(\ell)&0&128&192&224&240&248&252&254&255.
\end{array}
\]

\Lemma{inv_stage_bound_range}
For every legal inverse length, \(\code{inv_stage_bound}(\ell)\) lies between
0 and 24.
\Proof
Case analysis over the legal inverse lengths gives one of the values
\(16,\ldots,24\).

\Lemma{inv_stage_bound_pos}
For every legal inverse length, \(\code{inv_stage_bound}(\ell)\ge 0\).
\Proof
This is immediate from \code{inv_stage_bound_range}.

\Lemma{inv_stage_bound_lt31}
For every legal inverse length, \(\code{inv_stage_bound}(\ell)<31\).
\Proof
\code{inv_stage_bound_range} gives an upper bound of 24.

\Lemma{inv_stage_bound_ge16}
For every legal inverse length, \(\code{inv_stage_bound}(\ell)\ge 16\).
\Proof
Each case in the explicit definition is one of \(16,\ldots,24\).

\Lemma{inv_stage_next}
If \(\ell<256\) is a legal inverse length, then
\[
\code{inv_stage_bound}(2\ell)
=
\code{inv_stage_bound}(\ell)+1.
\]
\Proof
Doubling moves one step to the right in the inverse bound table.

\Lemma{inv_len_next_ok}
If \(\ell<256\) is a legal inverse length, then \(2\ell\) is legal.
\Proof
The legal values are powers of two from 1 through 256.  Doubling a value below
256 gives the next value in the list.

\Lemma{inv_zbase_range}
For every legal inverse length, \(\code{inv_zbase}(\ell)\) lies between 0 and
255.
\Proof
Immediate from the explicit inverse base table.

\Lemma{inv_read_range}
Assume a legal inverse length \(\ell<256\), \(0\le start<256\), and
\[
start=2\ell\,(zc-\code{inv_zbase}(\ell)).
\]
Then \(zc\in\Range{0}{254}\).
\Proof
The equation identifies \(zc-\code{inv_zbase}(\ell)\) with the completed block
count in the current inverse layer.  Since \(start\) is a valid block start and
\(\ell<256\), the layer consumes only indices below 255; checking the finite
base table gives \(zc\in\Range{0}{254}\).

\Lemma{inv_block_end_bound}
Under the same hypotheses,
\[
start+2\ell\le 256.
\]
\Proof
The start position is a valid block beginning in a layer whose block size is
\(2\ell\), and legal inverse lengths divide 256.

\Lemma{inv_stage_zbase_exit}
If an inverse layer exits with
\[
256=2\ell\,(zc-\code{inv_zbase}(\ell)),
\]
then \(zc=\code{inv_zbase}(2\ell)\).
\Proof
At layer exit, the number of consumed twiddle entries is the number of blocks
in the layer.  Adding that count to the old base gives the next base by the
explicit inverse table.

\Lemma{inv_middle_to_inner}
At block start, the inverse middle bound equals the inverse inner bound with
\(j=start\).
\Proof
No coefficient of the new block has been processed yet; unfolding the two
definitions leaves exactly the same processed prefix \(0\le i<start\).

\Lemma{inv_inner_current}
For \(j\) inside an inverse block, the current inputs \(j\) and \(j+\ell\) have
the old stage bound \(s\), and after advancing to \(j+1\) both have bound
\(s+1\).
\Proof
The inverse inner bound has the same shape as the forward one.  The current
pair is outside the completed subranges before the step and inside them after
the step.

\Lemma{inv_inner_weaken_unchanged}
For indices other than the current pair, the inverse inner bound can only stay
the same or increase when \(j\) advances to \(j+1\).
\Proof
Unfold the definition.  Advancing the loop extends the completed ranges by
exactly the two current positions.

\Lemma{inv_inner_exit_to_middle}
At inverse inner-loop exit, the inverse inner bound is the middle bound for
the next block start.
\Proof
At \(j=start+\ell\), the entire block \(start,\ldots,start+2\ell-1\) has been
processed.  This is exactly the prefix extension represented by
\(\code{inv_middle_bound}(s,start+2\ell,-)\).

\Lemma{inv_middle_exit_const}
If \(start\ge 256\), every in-range index has inverse middle bound \(s+1\).
\Proof
All indices \(0,\ldots,255\) lie below \(start\), so they are in the completed
prefix.

\Lemma{inv_middle_start_const}
At \(start=0\), the inverse middle bound is the constant bound \(s\).
\Proof
The completed prefix is empty.

\Lemma{inverse_spec_updateE}
The generated inverse update expression is extensionally equal to the compact
mathematical update
\[
(p[j\leftarrow p[j]+p[j+\ell]])[j+\ell\leftarrow z(p[j]-p[j+\ell])].
\]
\Proof
Since \(\ell>0\), the indices \(j\) and \(j+\ell\) are distinct.  Extensional
array equality and get-after-set case analysis reduce both sides to the same
value at every index.

\section{Inverse Butterfly and Twiddle Lemmas}

\Lemma{inverse_twiddle_product_bound_call_h}
For a valid inverse twiddle index \(zc\), if \(ai=\code{jzetas_inv}[zc]\) and
\(bi\) represents the field difference \(diff\), and the signed product bound
holds, then \code{__fqmul} returns a 16-bit word representing
\[
\code{zetas_inv}[zc]\cdot diff.
\]
\Proof
\code{fqmul_word_to_coeff_mul_bound_h} gives
\[
\operatorname{word}(res)=\operatorname{word}(ai)\operatorname{word}(bi)R^{-1}
\]
and a 16-bit result bound.  Substitute \(\operatorname{word}(bi)=diff\), and
use \code{jzetas_inv_get_cancel} to replace
\(\operatorname{word}(ai)R^{-1}\) by \(\code{zetas_inv}[zc]\).

\Lemma{inverse_twiddle_product_bound_call_ph}
The probability-one version of \code{inverse_twiddle_product_bound_call_h}.
\Proof
Combine the ordinary Hoare triple with \code{Fq.fqmul_ll}.

\Lemma{inverse_twiddle_value}
If a word \(t\) satisfies the Montgomery product equation with an inverse
twiddle and the field difference \(p[j]-p[j+\ell]\), then
\[
\operatorname{word}(t)=\code{zetas_inv}[zc]\,(p[j]-p[j+\ell]).
\]
\Proof
Substitute \(ai=\code{jzetas_inv}[zc]\) and use
\code{jzetas_inv_get_cancel}.  The remaining equality is associativity in
\(\Fq\).

\Lemma{inverse_butterfly_store}
If \(rp\sim p\), \(t_0=rp[j]\), \(coeff=rp[j+\ell]\), and \(t\) represents the
desired target value, then the stores
\[
rp[j]\leftarrow t_0+coeff,\qquad rp[j+\ell]\leftarrow t
\]
represent the inverse butterfly update
\[
p[j]\leftarrow p[j]+p[j+\ell],\qquad p[j+\ell]\leftarrow target.
\]
\Proof
The representation relation identifies \(t_0\) and \(coeff\) with
\(p[j]\) and \(p[j+\ell]\).  The lemma \code{word_to_coeff_add} converts
\(t_0+coeff\) into the field sum.  Two applications of
\code{NTT_Fq.poly_repr_set32} prove the representation after the two stores.

\Lemma{inverse_butterfly_store_bound}
Under a uniform \(s\)-bit bound and a 16-bit bound on \(t\), the inverse
butterfly stores produce a representation bounded by \(\max(s,16)+1\).
\Proof
The representation part is \code{inverse_butterfly_store}.  The bound for
\(t_0+coeff\) follows from \code{Fq.add_corr}; the bound for \(t\) is weakened
from 16 bits.  Untouched entries are weakened from \(s\) to the larger bound,
and the two store-bound lemmas finish the proof.

\Lemma{inverse_butterfly_store_bound_by}
This is the pointwise-bound version of the previous lemma.  The two current
inputs have bound \(a_s\); the target bound raises both updated positions to
\(a_s+1\).
\Proof
Use \code{inverse_butterfly_store} for representation and
\code{inverse_butterfly_bounds} for the new sum bound.  The freshly multiplied
word is 16-bit and therefore weakens to \(a_s+1\).  The two pointwise
get-after-set lemmas handle the stores.

\Lemma{inverse_butterfly_step_bound_by}
One full inverse butterfly step preserves representation and advances the
inverse inner-loop bound from \(j\) to \(j+1\).
\Proof
\code{inv_inner_current} supplies the bounds on the two inputs.
\code{inverse_butterfly_store_bound_by} performs the store proof.  The
untouched-lane hypothesis is supplied by \code{inv_inner_weaken_unchanged}.

\Lemma{inverse_twiddle_product_bound_from_inner}
Inside the inverse inner-loop invariant, the word difference
\[
rp[j]-rp[j+\ell]
\]
has bound \(\code{inv_stage_bound}(\ell)+1\).
\Proof
\code{inv_inner_current} gives the same stage bound for the two current input
words.  Applying \code{Fq.sub_corr} to those two bounded words gives one
additional bit for the difference.

\Lemma{inverse_twiddle_product_call_bound_from_inner}
Inside the inverse inner-loop invariant, the product of the inverse twiddle
and the word difference lies in the signed Montgomery reducer input range.
\Proof
The inverse twiddle is 16-bit by \code{jzetas_inv_bound16}.  The difference is
bounded by \code{inv_stage_bound} at length \(\ell\), plus one bit, by the previous lemma; for
\(\ell<256\) this is at most 24.  Weakening gives a 24-bit bound, and
\code{fqmul_product_bound_16_24} supplies the reducer range.

\Lemma{inverse_butterfly_step_full}
If \(t\) satisfies the Montgomery product equation for the inverse twiddle and
the field difference, then the two stores represent the abstract inverse
butterfly update.
\Proof
Use \code{inverse_twiddle_value} to convert the Montgomery product equation
into the abstract product with \(\code{zetas_inv}[zc]\).  Then apply
\code{inverse_butterfly_store}.

\Lemma{inverse_butterfly_repr}
This is another representation-only form of the inverse butterfly step, using
the same hypotheses as \code{inverse_butterfly_step_full}.
\Proof
The proof is identical in structure: derive the abstract value of \(t\) by
\code{inverse_twiddle_value}, then call \code{inverse_butterfly_store}.

\section{Final Inverse Scaling Lemmas}

The final inverse loop is tracked by the ghost array
\[
\code{final_mont_array}(p,j)[i]
=
\begin{cases}
p[i]R,&0\le i<j,\\
p[i],&\text{otherwise}.
\end{cases}
\]
The matching bound function is
\[
\code{final_bound}(j,i)
=
\begin{cases}
16,&0\le i<j,\\
24,&\text{otherwise}.
\end{cases}
\]

\Lemma{final_mont_array_start}
\[
\code{final_mont_array}(p,0)=p.
\]
\Proof
The prefix \(0\le i<0\) is empty, so the initializer returns \(p[i]\) at every
index.

\Lemma{final_mont_array_exit}
\[
\code{final_mont_array}(p,256)=\code{NTT_Fq.array256_mont}(p).
\]
\Proof
Every array index lies in the prefix \(0\le i<256\).  Hence the initializer
returns \(p[i]R\) at each index, which is exactly \code{array256_mont} by
\code{array256_mont_get}.

\Lemma{final_mont_array_get_raw}
If \(0\le j<256\), then the \(j\)-th entry of
\(\code{final_mont_array}(p,j)\) is still \(p[j]\).
\Proof
The condition for Montgomery scaling is \(0\le i<j\).  At \(i=j\), this
condition is false, so the raw entry \(p[j]\) is returned.

\Lemma{final_mont_array_setE}
For \(0\le j<256\), storing \(p[j]\cdot\code{scale255}\cdot R\) at index \(j\)
of \(\code{final_mont_array}(p,j)\) yields
\[
\code{final_mont_array}(p[j\leftarrow p[j]\cdot\code{scale255}],j+1).
\]
\Proof
Use extensional equality.  At \(i=j\), both sides are the newly scaled value.
For \(i<j\), both sides are already Montgomery-scaled.  For \(i>j\), both
sides are still raw entries.  The get-after-set rules separate these cases.

\Lemma{final_bound_start}
\[
\code{final_bound}(0,i)=24.
\]
\Proof
No index lies in the processed prefix \(0\le i<0\).

\Lemma{final_bound_next_unchanged}
If \(k\ne j\), then the bound at \(k\) under \(\code{final_bound}(j,-)\) is no
larger than the bound under \(\code{final_bound}(j+1,-)\).
\Proof
Advancing \(j\) changes only the status of index \(j\), moving it from the
24-bit unprocessed region to the 16-bit processed region.  Every other index
keeps the same status, or is already in the processed prefix.

\Lemma{inverse_final_product_bound_call_h}
For the special table entry \(ai=\code{jzetas_inv}[255]\), if \(bi\) represents
\(coeff\) and the signed product bound holds, then \code{__fqmul} returns a
16-bit word representing
\[
coeff\cdot \code{scale255}\cdot R.
\]
\Proof
\code{fqmul_word_to_coeff_mul_bound_h} gives the Montgomery product
\(\operatorname{word}(ai)\operatorname{word}(bi)R^{-1}\) and the 16-bit bound.
The hypothesis identifies \(\operatorname{word}(bi)=coeff\), and
\code{jzetas_inv_255_scale_cancel} rewrites
\(\operatorname{word}(ai)R^{-1}\) as \(\code{scale255}R\).

\Lemma{inverse_final_product_bound_call_ph}
The probability-one version of \code{inverse_final_product_bound_call_h}.
\Proof
Use losslessness of \code{__fqmul} and the ordinary Hoare triple.

\Lemma{inverse_final_product_bound_from_inv}
Under the final-loop bound invariant, the product of
\(\code{jzetas_inv}[255]\) and the current unprocessed coefficient word lies in
the signed Montgomery reducer input range.
\Proof
\code{jzetas_inv_bound16} gives the 16-bit bound for the scale table entry.
The current unprocessed coefficient has 24-bit bound by \code{final_bound}.
Apply \code{fqmul_product_bound_16_24}.

\section{Wrapper and Core pRHL Equivalences}

\Lemma{poly_ntt_jazz_ref}
The public extracted wrapper \code{poly_ntt_jazz} is equivalent to the core
procedure \code{_poly_ntt} under equality of the input buffer and equality of
the result.
\Proof
The wrapper contains only the core call.  Inlining all procedures reduces both
sides to the same program, and pRHL simulation proves equality of results.

\Lemma{poly_invntt_jazz_ref}
The public extracted wrapper \code{poly_invntt_jazz} is equivalent to the core
procedure \code{_poly_invntt}.
\Proof
As in the forward case, inlining exposes a direct call to the core routine, and
the two sides simulate identically.

\Lemma{poly_ntt_core_ref}
The extracted forward core refines the imperative EasyCrypt forward NTT:
\[
\begin{array}{l}
\code{Hpoly_extract.M._poly_ntt}
\sim
\code{NTT_Fq.NTT.ntt},\\[0.3em]
\code{poly_repr_bound}(rp_1,r_2,16)\land
\code{zetas}_2=\code{NTT_Fq.zetas}\\
\qquad\Longrightarrow
\code{poly_repr_bound}(res_1,res_2,24).
\end{array}
\]
\Proof
The proof is a direct three-loop pRHL simulation.  The outer invariant aligns
the layer length, twiddle counter, table pointers, representation relation,
and a uniform pointwise bound
\(\code{fwd_stage_bound}(\ell)\).  The middle invariant refines the bound to
\code{fwd_middle_bound}, distinguishing completed blocks from untouched
blocks.  The inner invariant refines it further to \code{fwd_inner_bound},
distinguishing completed butterflies inside the current block.

At each inner iteration, \code{forward_twiddle_product_bound_call_ph} verifies
the call to \code{__fqmul}; its product-range precondition is supplied by
\code{forward_twiddle_product_bound_from_inner}.  The subsequent stores are
verified by \code{forward_butterfly_step_bound_by}.  Inner-loop exit uses
\code{fwd_inner_exit_to_middle}; middle-loop exit uses
\code{fwd_middle_exit_const}; outer-loop advancement uses
\code{fwd_stage_next}, \code{fwd_len_next_ok}, and \code{fwd_zbase_next}.  At
outer-loop exit, \(\ell=0\), so the bound is 24, giving the stated
postcondition.

\Lemma{poly_invntt_core_ref}
The extracted inverse core refines the imperative EasyCrypt inverse NTT:
\[
\begin{array}{l}
\code{Hpoly_extract.M._poly_invntt}
\sim
\code{NTT_Fq.NTT.invntt},\\[0.3em]
\code{poly_repr_bound}(rp_1,r_2,16)\land
\code{zetas_inv}_2=\code{NTT_Fq.zetas_inv}\\
\qquad\Longrightarrow
\code{poly_repr_bound}(res_1,\code{NTT_Fq.array256_mont}(res_2),16).
\end{array}
\]
\Proof
The inverse proof has two phases.  First, weakest-precondition reasoning
encounters the syntactically final scaling loop.  Its invariant represents the
left buffer as \code{final_mont_array} of the right buffer and uses
\code{final_bound} to say that the processed prefix is 16-bit while the
unprocessed suffix is 24-bit.  Each iteration is justified by
\code{inverse_final_product_bound_call_ph} and
\code{inverse_final_product_bound_from_inv}.  The lemmas
\code{final_mont_array_setE}, \code{final_bound_next_unchanged}, and
\code{final_mont_array_exit} carry the loop from an empty Montgomery prefix to
\code{array256_mont}.

Second, the proof establishes the eight inverse butterfly layers with the same
three-loop structure as the forward proof, using \code{inv_stage_bound},
\code{inv_middle_bound}, and \code{inv_inner_bound}.  The multiplication call
inside a butterfly is handled by \code{inverse_twiddle_product_bound_call_ph};
the product-range side condition comes from
\code{inverse_twiddle_product_call_bound_from_inner}.  The stores are handled
by \code{inverse_butterfly_step_bound_by}, and \code{inverse_spec_updateE}
aligns the right-hand array update with the generated program expression.
Layer advancement uses \code{inv_stage_next}, \code{inv_len_next_ok}, and
\code{inv_stage_zbase_exit}.  After the butterfly layers, the output is
represented with a 24-bit bound; this initializes the final scaling invariant.
The final loop then reduces every entry back to 16 bits and produces the
Montgomery-scaled representation required by the theorem.

\section{Conclusion}

Every lemma in \code{RefJasminNTT.ec} serves one of four mathematical roles:
\begin{enumerate}
\item interpreting machine words as field elements without overflow;
\item cancelling Montgomery factors in concrete twiddle tables;
\item maintaining pointwise word-size bounds through butterfly loops;
\item assembling those local facts into direct pRHL simulations of the
      extracted forward and inverse NTT cores.
\end{enumerate}
The final result is a checked refinement layer from Jasmin extraction to the
imperative EasyCrypt NTT model, with the inverse postcondition strengthened to
the mathematically correct Montgomery-scaled representation.

\end{document}
